% DRAFT — Adaptive Learning Systems section (Agent E/F/G output)
% Review before merging into literature_review.tex
% Pipeline: paper cards A–C → bib_mapping.csv → chronological B2 prose + expanded table
% Do NOT merge Emotion-Aware Adaptation section from this file.

\section{Adaptive Learning Systems}\label{sec:adaptive_learning_systems}

While the previous section focused on recognising learner states through facial behaviour, adaptive learning systems focus on deciding how instruction should respond to learner needs. These systems are digital tools that change how they teach depending on the individual learner. Instead of presenting the same material in the same order to every student, they adjust things like the content, the pace, the difficulty of the tasks, the feedback given after each answer, hints, and presentation format \cite{strielkowski_span_2025}. A common way to describe this behaviour is as a closed loop: the system collects data from the learner, uses that data to judge how the learner is progressing, and then decides what to show next \cite{gligorea_adaptive_2023}. Reviews of intelligent tutoring systems stress that adaptation is not limited to picking the next exercise: the tutor model may also schedule scaffolding, remediation, motivational messages, interface changes, and assessment pacing \cite{zerkouk_comprehensive_2025}. The goal is to give each learner tasks that match their current level, so that they can keep improving without feeling bored or overwhelmed \cite{zerkouk_comprehensive_2025}.

Early adaptive systems focused on explicit rules written by experts. These were called intelligent tutoring systems (ITS), and examples such as GUIDON and the LISP Tutor were built in the 1980s around hand-crafted if--then rules that decided when to give a hint, when to move on, and when to mark a concept as learned \cite{zerkouk_comprehensive_2025}. The strength of these systems was that their decisions were transparent and easy to explain, but the rules had to be prepared for every topic, which made them hard to extend beyond a single subject. Rule-based logic has not disappeared, though. Sayed et al., for example, used a small set of threshold rules based on a student's penalised grade to move an exercise up or down in difficulty inside an online platform for third-grade mathematics \cite{sayed_ai-based_2023}.

Later work introduced student models that tried to estimate what each learner already knew. Instead of relying only on teacher-written rules, these systems built a picture of the learner from their interaction history and updated it after every answer \cite{zerkouk_comprehensive_2025}. This line of research eventually led to knowledge tracing methods, in which the system predicts how likely the learner is to answer the next question correctly based on their earlier answers. Pardos and colleagues showed that year-long patterns of boredom, engaged concentration, confusion, and gaming-the-system behaviour in ASSISTments logs correlate with end-of-year mathematics test scores \cite{Pardos2014AffectiveStatesStateTests}. Kuling and Zitnik built a recent exercise recommender that tries to pick questions just above the learner's current level \cite{kuling_ken_2025}. Educational psychologists call this range the Zone of Proximal Development: the set of tasks that are challenging but still reachable. Their system ranks candidate questions by predicted difficulty and then chooses items around the 66th percentile, so that the next exercise is hard but not impossible \cite{kuling_ken_2025}.

Over time, researchers began to treat task selection as a decision problem that could be learned from experience. Reinforcement learning (RL) then became a common tool for this. In RL, the system tries different actions, observes the results, and slowly learns which actions lead to better outcomes. Azoulay et al. compared several adaptation methods within the same simulated environment. The study tested different reinforcement learning algorithms, a bandit-style method, and a Bayesian approach that keeps a probability estimate of each learner's ability \cite{azoulay_adaptive_2021}. The Bayesian method reached about 89 to 93 percent of the performance of an ideal algorithm that already knew each student's real ability, while the reinforcement learning methods reached lower but still useful levels \cite{azoulay_adaptive_2021}. The study has become an important reference because it placed very different adaptation families on equal ground, and later work built on its baselines \cite{perez_emotions_2024}.

More recent studies have combined knowledge tracing with reinforcement learning. Fu proposed RL-DKT, a system that first uses a deep knowledge tracing model to estimate the learner's current knowledge state, and then uses a reinforcement learning agent to choose the next task from that estimate \cite{fu_integrating_2025}. Fu tested the method on real learner logs from well-known datasets such as ASSISTments, KDD Cup 2010, and Cognitive Tutor, and reported AUC of 0.874, F1 of 0.812, average completion time of 34.5 minutes, and dropout of 3.2\% \cite{fu_integrating_2025}. Sayed et al. followed a different mix, using a Deep Q-Network to choose how to present the learning content while still using simple rules to move between difficulty levels, and reported learning gains in a small pilot with grade-three mathematics students \cite{sayed_ai-based_2023}. Other studies have used reinforcement learning mainly to decide when and how to give feedback or hints during a session \cite{deshmukh_developing_2025,guzman_developing_2025}. Guzman and Cruz-Mercado, for example, trained a Proximal Policy Optimisation agent that selects problems, hints, and instructional review, and reported normalised learning gains of 0.72 versus 0.58 for a rule-based tutor in a study of ninety introductory Python learners \cite{guzman_developing_2025}. Taken together, these works show that adaptation is increasingly treated as a step-by-step decision problem over tasks, hints, feedback, pacing, or review, in which each chosen action builds on what the learner has done so far.

From 2024, several systems began to close the loop from sensed affect to more than difficulty alone. P�rez et al. used valence and arousal from facial expressions as implicit feedback for a Q-learning difficulty controller \cite{perez_emotions_2024}. Govea and colleagues proposed a deeper variant, in which a Proximal Policy Optimisation agent adapted pace, difficulty and activity type based on a multimodal state that combined face, voice and wrist biometrics \cite{govea_implementation_2024}. Gong described a similar multimodal architecture for e-learning, built around a Deep Q-Network that adjusts difficulty through Item Response Theory, pacing through a Markov-modulated policy and content format between visual, auditory and textual, as well as hints and motivational feedback \cite{gong_design_2025}. Kong and colleagues combined facial expression, voice, heart-rate variability, electrodermal activity and optional electroencephalography into a single joint state of cognitive load and achievement emotions, and a contextual bandit with a Deep Q-Network back end selected interventions such as worked examples, scaffolds, pacing changes or motivational feedback \cite{kong_real-time_2025}. Their randomised evaluation with two hundred and seventy-one learners, in both laboratory and classroom settings, is one of the strongest current examples of a fully working emotion-aware adaptive tutor \cite{kong_real-time_2025}. A scoping review by Fern�ndez-Herrero, also from 2024, confirmed that scaffolding, content adaptation, difficulty, and feedback remain the dominant pedagogical responses when emotion is sensed \cite{fernandez-herrero_evaluating_2024}. These systems are the closest predecessors to an emotion-aware reinforcement-learning tutor, but they either rely on simulation, on static expression corpora, or on sensor stacks that are difficult to deploy in ordinary online classrooms. How these emotion-aware ideas have been put into practice, and how they connect to the adaptive methods described here, is the focus of the next section.

Very recent work extends adaptive tutoring toward generative AI and standardised evaluation. Meng and Yang describe a multi-agent tutor in which Proximal Policy Optimisation selects interventions while neural knowledge tracing and an engagement model track mastery and participation \cite{meng_self-evolving_2025}. Gutierrez, Villegas-Ch and Govea moved to the language-model era with a conversational assistant that adapts feedback, difficulty and communication tone from chat messages \cite{gutierrez_development_2025}. A recent testbed called TutorGym was built to make step-level benchmarking easier and more comparable across studies \cite{weitekamp_tutorgym_2025}.

Evaluating adaptive systems is difficult, because running long classroom studies with many learners is expensive and slow. To deal with this, many studies first test their methods in simulators or on stored learner logs before any live deployment. Azoulay et al. used an artificial environment with 10{,}000 simulated runs to compare adaptation algorithms against a theoretical upper bound \cite{azoulay_adaptive_2021}. P�rez et al. followed a similar style, running 1{,}000 simulations over 200 tasks to compare their approach against earlier reinforcement learning baselines \cite{perez_emotions_2024}. Studies such as those by Fu and by Kuling and Zitnik instead use offline learner datasets and report standard metrics like prediction accuracy and ranking scores on real interaction logs \cite{fu_integrating_2025,kuling_ken_2025}. When live testing is possible, it still tends to be small: Kuling and Zitnik ran a thirteen-week classroom deployment with 38 graduate students \cite{kuling_ken_2025}, and Sayed et al. reported a two-month pilot with 26 grade-three learners \cite{sayed_ai-based_2023}. These practical evaluation approaches matter because they let researchers compare methods and build confidence in them before committing to a full classroom study.

Despite the variety of methods, the main strengths of adaptive learning systems are fairly consistent across the literature. They can keep each learner close to a suitable level of challenge, deliver feedback faster than a human teacher could in a large class, and support thousands of learners at once \cite{gligorea_adaptive_2023}. Large commercial platforms such as Knewton, Carnegie Learning's MATHia, DreamBox, and Smart Sparrow already use these ideas to deliver personalised mathematics and other subjects at scale \cite{strielkowski_span_2025}. Reviews of the field agree that these systems offer clear benefits for engagement and performance when content is well structured and adaptation signals are rich enough \cite{gligorea_adaptive_2023,zerkouk_comprehensive_2025}.

Despite these strengths, the signals that drive the adaptation are still quite narrow. In most of the studies discussed above, the system decides what to show next based only on what the learner did: whether the answer was correct, how long the answer took, how many hints were used, and which items were clicked \cite{azoulay_adaptive_2021,fu_integrating_2025,kuling_ken_2025}. Even when tutors adjust pacing, feedback, or presentation format, the richest multimodal designs still depend on sensors beyond an ordinary webcam \cite{govea_implementation_2024,kong_real-time_2025,gong_design_2025}. How the learner feels during the task, whether they are confused, bored, or frustrated, rarely enters the decision. Azoulay et al. themselves acknowledge this gap, noting that their simulation could check whether the algorithm picked the right difficulty, but could not reflect the learners' satisfaction, motivation, or real learning gain \cite{azoulay_adaptive_2021}. A recent comprehensive review of AI-based intelligent tutoring systems lists affective computing as one of the main open challenges for the field, adding that these systems are still less proficient than human tutors at interpreting students' emotional states \cite{zerkouk_comprehensive_2025}. A recent scoping review reports that more researchers are now trying to add emotion as an additional input channel in intelligent tutoring systems \cite{fernandez-herrero_evaluating_2024}. How these emotion-aware ideas have been put into practice, and how they connect to the adaptive methods described here, is the focus of the next section. To summarise the most relevant adaptive learning studies discussed in this section, Table~\ref{tab:adaptive_systems_comparison} compares their adaptation targets, input signals, decision methods, evaluation settings, main reported results, and relevance to the present thesis.

\clearpage
\begin{landscape}
\thispagestyle{empty}
\small
\setlength{\tabcolsep}{3pt}
\renewcommand{\arraystretch}{1.1}

\begin{longtable}{>{\RaggedRight\arraybackslash}p{2.5cm}
                  >{\RaggedRight\arraybackslash}p{2.1cm}
                  >{\RaggedRight\arraybackslash}p{2.2cm}
                  >{\RaggedRight\arraybackslash}p{2.7cm}
                  >{\RaggedRight\arraybackslash}p{2.0cm}
                  >{\RaggedRight\arraybackslash}p{2.8cm}
                  >{\RaggedRight\arraybackslash}p{3.2cm}
                  >{\RaggedRight\arraybackslash}p{3.4cm}}
\caption{Comparison of key adaptive learning studies relevant to this thesis}
\label{tab:adaptive_systems_comparison} \\

\toprule
\textbf{Paper} & \textbf{Adaptation target} & \textbf{Input signals} & \textbf{Decision method} & \textbf{Evaluation type} & \textbf{Main quantitative results} & \textbf{Main strength / limitation} & \textbf{Relevance to this thesis} \\
\midrule
\endfirsthead

\toprule
\textbf{Paper} & \textbf{Adaptation target} & \textbf{Input signals} & \textbf{Decision method} & \textbf{Evaluation type} & \textbf{Main quantitative results} & \textbf{Main strength / limitation} & \textbf{Relevance to this thesis} \\
\midrule
\endhead

\bottomrule
\endfoot

\makecell[l]{Azoulay et al.\\ (2021)\\ \cite{azoulay_adaptive_2021}} &
Next-task / difficulty selection &
Simulated learner performance, success/failure history &
Comparison of Q-learning, TD, VL, DVRL, Gittins, Rasch+Elo, Bayesian inference &
Simulation &
10{,}000 runs; Bayesian method reached 89--93\% of oracle utility &
Canonical multi-method benchmark; simulation only, no affect or real learners &
Establishes RL/IRT baselines for next-task selection without emotion in state \\

\makecell[l]{Sayed et al.\\ (2023)\\ \cite{sayed_ai-based_2023}} &
Content presentation, exercise navigation, feedback &
Learner model (cognitive, behavioural, affective cues), VARK style &
DQN + threshold rules for difficulty &
Classroom pilot (grade 3, $n=26$, 2 months) &
Post-test gains; largest gains for initially low performers &
Hybrid rules + DQN in deployed e-learning; affect channel not FER-only &
Shows hybrid RL deployment in schools; contrast for full RL + webcam FER policy \\

\makecell[l]{Fu\\ (2025)\\ \cite{fu_integrating_2025}} &
Personalized learning path / next-task sequencing &
Interaction logs, DKT-estimated knowledge state &
RL-DKT: REINFORCE + dynamic knowledge tracing &
Offline datasets (ASSISTments, KDD, Cognitive Tutor) &
AUC 0.874; F1 0.812; avg.\ completion 34.5 min; dropout 3.2\%; path efficiency +12.5\% &
Strong KT+RL sequencing on real logs; no facial affect in state &
Representative performance-only RL sequencing the thesis extends with FER \\

\makecell[l]{Kuling and\\ Zitnik\\ (2025)\\ \cite{kuling_ken_2025}} &
Adaptive exercise / next-question recommendation &
Interaction history, correctness, text embeddings &
KUL-Rec dual-memory Hebbian recommender &
10 offline datasets + classroom ($n=38$, 13 weeks) &
nDCG 0.316 vs 0.265 baseline; Recall@10 0.305 vs 0.211; helpfulness $p<.05$ &
Strong ZPD-targeted recommendation; not RL; no emotion &
Shows next-item adaptation from sparse traces without affect sensing \\

\makecell[l]{Guzman and\\ Cruz-Mercado\\ (2025)\\ \cite{guzman_developing_2025}} &
Problem selection, hints, instructional review &
DKVMN knowledge state; synthetic then real traces &
PPO with multi-faceted reward (performance, efficiency, retention) &
Simulation (10{,}000 synthetic) + human study ($N=90$) &
Normalised learning gains 0.72 vs 0.58 (rule-based) vs 0.45 (static), $p<0.01$ &
End-to-end PPO tutor with human subjects; no facial affect &
Direct PPO comparator for emotion-augmented policy design \\

\makecell[l]{Deshmukh and\\ Sen\\ (2025)\\ \cite{deshmukh_developing_2025}} &
Feedback timing and type (hints, explanations, encouragement) &
Simulated correctness, time on task, hint requests &
Tabular Q-learning &
Simulation vs rule-based and random feedback &
Test accuracy 85\% vs 73\% (rule-based) vs 66\% (random); satisfaction 4.4/5 vs 3.8 &
Simple RL feedback loop; short study; no KT or affect &
Illustrates RL over feedback actions adjacent to thesis intervention set \\

\makecell[l]{P\'erez et al.\\ (2024)\\ \cite{perez_emotions_2024}} &
Difficulty adaptation / next task &
Correctness, BKT mastery, valence/arousal (circumplex) &
Q-learning with emotion as implicit feedback in reward and exploration &
Simulation (1{,}000 runs, 200 tasks) &
Learning gain $\approx$0.49 vs 0.28 (Q-learning) at 200 tasks &
Closest prior: affect in RL loop; simulation and static expression corpus &
Very highly relevant: emotion-aware difficulty RL the thesis generalises to multi-action FER \\

\makecell[l]{Govea et al.\\ (2024)\\ \cite{govea_implementation_2024}} &
Pace, difficulty, activity type, presentation &
Face, voice, wrist biometrics &
CNN--RNN + Proximal Policy Optimization &
Hybrid classroom; multimodal sensing &
Emotion detection 72.4\% $\rightarrow$ 89.3\%; personalization 70.2\% $\rightarrow$ 90.1\% &
Multimodal affect closed loop; adaptation evidence still limited &
Relevant multimodal emotion-to-policy precedent; thesis uses webcam-only sensing \\

\makecell[l]{Gong\\ (2025)\\ \cite{gong_design_2025}} &
Difficulty (IRT), pacing, content format, hints, motivational feedback &
Face (CNN), speech (MFCC+BiLSTM), text (BERT) &
Deep Q-Network adaptive delivery engine &
Controlled e-learning pilot &
Fusion accuracy 0.91; F1 0.90; engagement $\approx$0.92 vs 0.60 baseline &
Multimodal emotion-aware RL over several pedagogical levers; not FER-only &
Very highly relevant multi-action RL design; thesis narrows to facial input \\

\makecell[l]{Kong et al.\\ (2025)\\ \cite{kong_real-time_2025}} &
Worked examples, scaffolds, pacing, motivational feedback, content/interface &
Face, voice, HRV, EDA, optional EEG, interaction logs &
Contextual bandit + Deep Q-Network back end &
RCT ($N=271$); lab + classroom &
Weighted F1 $\approx$0.85; mid-ability post-test gain $\approx$+15 pp; engagement 42\% vs 29\% &
Strongest field evaluation of multi-intervention affect-aware tutor; heavy sensing &
Sets breadth benchmark for intervention set; thesis targets lighter FER+RL stack \\

\makecell[l]{Meng and\\ Yang\\ (2025)\\ \cite{meng_self-evolving_2025}} &
Intervention selection (hints, scaffolding, presentation) &
SAINT+ mastery, behavioural engagement, LLM responses &
PPO multi-agent tutor (RLTS) &
Simulated replay of historical logs &
+28.6\% adaptability; $-$31.2\% recurring errors; $-$24.8\% dropout vs static tutors &
GenAI + KT + RL stack; behavioural not facial engagement &
Contemporary PPO tutor without webcam FER in state \\

\makecell[l]{Weitekamp\\ et al.\\ (2025)\\ \cite{weitekamp_tutorgym_2025}} &
Step-level tutor/student decisions (benchmark) &
ITS step state (223 domains: CTAT, OATutor, etc.) &
Standard API for LLMs, RL, cognitive models &
Benchmark evaluation across existing ITS interfaces &
LLMs near chance on incorrect-action labelling; next-step accuracy $\approx$52--70\% &
Infrastructure for comparable step-level evaluation; not an adaptation algorithm &
Supports simulation-first, step-level evaluation rationale for thesis RL module \\

\end{longtable}
\end{landscape}
\clearpage
